\documentclass[11pt,a4paper]{article}
% Common preamble for the hint sheets (\input from each file)
\usepackage[utf8]{inputenc}
\usepackage[T1]{fontenc}
\usepackage{lmodern}
\usepackage{amsmath,amssymb,amsthm,mathtools}
\usepackage[a4paper,margin=2.3cm]{geometry}
\usepackage{enumitem}
\usepackage{xcolor}
\usepackage[hidelinks]{hyperref}
\renewcommand{\leq}{\leqslant}
\renewcommand{\le}{\leqslant}
\renewcommand{\geq}{\geqslant}
\renewcommand{\ge}{\geqslant}
\setlength{\parindent}{0pt}
\setlength{\parskip}{0.35em}
% \hintstitle{Sheet title}{Class work / Homework}
\newcommand{\hintstitle}[2]{\begin{center}{\LARGE\bfseries Hints: #1\par}\vspace{0.5em}{\large #2\par}\end{center}
\noindent\textit{The hints to each problem are cumulative. Read one, try again, and open the next only if you are still stuck.}\par\medskip\hrule\medskip}
\newcommand{\problem}[1]{\par\bigskip{\large\bfseries Problem #1}\par\nopagebreak\smallskip}
\newcommand{\subpart}[1]{\par\smallskip\textbf{(#1)}\par\nopagebreak}
\newcommand{\hint}[1]{\par\smallskip\noindent\textcolor{black!70}{\textbf{Hint #1.}}\enspace}
\newcommand{\note}[1]{\par\smallskip\noindent\textit{Note.} #1\par}
\providecommand{\lt}{<}
\providecommand{\gt}{>}
\providecommand{\tightlist}{\setlength{\itemsep}{0pt}\setlength{\parskip}{0pt}}
\emergencystretch=3em

\begin{document}
\hintstitle{Limits -- 2. Recurrences solvable in closed form}{Homework}

\problem{1}
\hint{1} Let $f(x)=(x^2-9)/8$. Its fixed points are $9$ and $-1$; note $f'(9)=\frac94$ (repelling) and $f'(-1)=-\frac14$ (attracting). Also $f$ is even and $f(x)\ge-\frac98$, so only $|x_1|$ matters and $x_n\ge-\frac98$ for $n\ge2$.
\hint{2} Use $f(x)-9=\frac{(x-9)(x+9)}8$: if $|x_1|>9$ then $x_2>9$ and $x_n-9$ grows at least geometrically, so $x_n\to+\infty$; if $|x_1|=9$ then $x_n=9$ for $n\ge2$.
\hint{3} If $|x_1|<9$ then $x_2\in[-\frac98,9)$. Use $f(x)+1=\frac{(x+1)(x-1)}8$: for $x\in[-\frac98,9)$ we have $|x-1|<8$, so $d_n=|x_n+1|$ decreases; then all later $x_n$ lie in $[-\frac98,-1+d_2]$, where $|x-1|\le 8q$ for a fixed $q<1$, so $d_n\to0$ geometrically.
\hint{4} Answer: $\lim x_n=-1$ if $|x_1|<9$, $=9$ if $|x_1|=9$, $=+\infty$ if $|x_1|>9$.

\problem{2}
\hint{1} Uniqueness and existence: $g(x)=x-\varepsilon\sin x$ is continuous with $g'(x)=1-\varepsilon\cos x\ge1-|\varepsilon|>0$, and $g(x)\to\pm\infty$ as $x\to\pm\infty$.
\hint{2} Rewrite the equation as $x=t+\varepsilon\sin x$ and iterate $x_{n+1}=t+\varepsilon\sin x_n$ from any $x_0$. Since $|\sin u-\sin v|\le|u-v|$, the map is a contraction with constant $|\varepsilon|<1$: $|x_{n+1}-x|\le|\varepsilon|\,|x_n-x|$, so $x_n\to x$ geometrically (this is Kepler's equation).

\problem{3}
\hint{1} Solve the recurrence: the characteristic roots are $2$ and $\frac12$; find $F(n)$ in closed form.
\hint{2} $F(n)=2^n-2^{-n}$, so with $y_n=2^{2^n}$ the terms are $\dfrac1{y_n-1/y_n}=\dfrac{y_n}{y_n^2-1}$, and $y_{n+1}=y_n^2$.
\hint{3} $\dfrac{y}{y^2-1}=\dfrac1{y-1}-\dfrac1{y^2-1}$, so the series telescopes: $\sum_{n\ge0}\frac1{F(2^n)}=\frac1{y_0-1}=1$. Rational.

\problem{4}
\hint{1} Compute: $1,2,6,22,90,394,\dots$ These are the (large) Schr\"oder numbers, which satisfy the integral recurrence $s_n=s_{n-1}+\sum_{k=0}^{n-1}s_ks_{n-1-k}$, i.e.\ their generating function $S(x)$ satisfies $S=1+xS+xS^2$. The task is to prove $a_n=s_n$.
\hint{2} Translate the given recurrence into an equation for $f(x)=\sum a_nx^n$: multiplying by $x^{n+2}$ and summing gives the first-order linear ODE
\[x(1-6x+x^2)f'(x)+(1-3x)f(x)=1+x .\]
A formal power series solution of this ODE is unique (the coefficient of $x^m$ determines $a_m$), so it suffices to check that $S$ satisfies it.
\hint{3} From $S=1+xS+xS^2$ get $S(x)=\dfrac{1-x-\sqrt{1-6x+x^2}}{2x}$; differentiate and verify the ODE. Then $a_n=s_n$, which are integers by the convolution recurrence.

\problem{5}
\hint{1} Use the generating function $f(x)=\sum a_nx^n$ (by induction $0<a_n\le1$, so it converges for $|x|<1$). The required limit is $f(\frac12)$.
\hint{2} The recurrence says $(n+1)a_{n+1}=\sum_{k=0}^n a_k\cdot\frac1{n-k+2}$, i.e.\ $f'(x)=f(x)g(x)$ with $g(x)=\sum_{m\ge0}\frac{x^m}{m+2}=\dfrac{-\ln(1-x)-x}{x^2}$.
\hint{3} An antiderivative of $g$ vanishing at $0$ is $1+\frac{(1-x)\ln(1-x)}{x}$, so $f(x)=e\,(1-x)^{(1-x)/x}$ and the limit is $f(\frac12)=\dfrac e2$.

\problem{6}
\hint{1} Decompose into partial fractions: $\dfrac1{(4k+1)(4k+2)(4k+3)(4k+4)}=\dfrac16\Bigl(\dfrac1{4k+1}-\dfrac3{4k+2}+\dfrac3{4k+3}-\dfrac1{4k+4}\Bigr)$.
\hint{2} Write $\frac1{4k+j}=\int_0^1x^{4k+j-1}\,dx$ and sum the geometric series (justify swapping, e.g.\ by monotone convergence on $[0,1)$): the sum is $\frac16\int_0^1\frac{(1-x)^3}{1-x^4}\,dx=\frac16\int_0^1\frac{(1-x)^2}{(1+x)(1+x^2)}\,dx$.
\hint{3} $\frac{(1-x)^2}{(1+x)(1+x^2)}=\frac2{1+x}-\frac{x+1}{1+x^2}$, giving the answer $\dfrac{\ln2}4-\dfrac\pi{24}$.

\problem{7}
\hint{1} The map $x\mapsto x^2-2$ is $t+t^{-1}\mapsto t^2+t^{-2}$. Write $x_1=\sqrt5=t+t^{-1}$.
\hint{2} $t=\frac{1+\sqrt5}2$, so $x_n=t^{2^{n-1}}+t^{-2^{n-1}}$, and the product telescopes: $(t-t^{-1})x_1x_2\cdots x_n=t^{2^n}-t^{-2^n}$.
\hint{3} Hence the ratio equals $\dfrac{1}{t-t^{-1}}\cdot\dfrac{t^{2^n}-t^{-2^n}}{t^{2^n}+t^{-2^n}}\to\dfrac1{t-t^{-1}}=1$.

\problem{8}
\hint{1} The sum is the "odd part" of a binomial expansion: $S_n=\sum_k\binom n{2k+1}5^k=\dfrac{(1+\sqrt5)^n-(1-\sqrt5)^n}{2\sqrt5}$.
\hint{2} By Theorem 2, $S_n$ satisfies $S_{n+1}=2S_n+4S_{n-1}$ ($1\pm\sqrt5$ are the roots of $x^2-2x-4$), with $S_1=1$, $S_2=2$. Prove $2^{n-1}\mid S_n$ by strong induction. (Equivalently, $S_n=2^{n-1}F_n$.)

\problem{9}
\hint{1} The number of ways is the constant term of $(x^{-1}+1+x)^{p-1}$. Choosing first which $2k$ minutes are jumps and then which $k$ of them go to the right, $f(p)=\sum_k\binom{p-1}{2k}\binom{2k}k$ (check: $f(3)=3$, $f(5)=19$, $f(7)=141$). For $p=2$: $f(2)=1$.
\hint{2} Modulo $p$ one has $\binom{p-1}{j}\equiv(-1)^j$, so $f(p)\equiv\sum_{k=0}^{(p-1)/2}\binom{2k}k\pmod p$ for odd $p$. With $m=\frac{p-1}2\equiv-\frac12$ one has $\binom{2k}k=(-4)^k\binom{-1/2}{k}\equiv(-4)^k\binom mk$ (all denominators are $k!$ with $k<p$).
\hint{3} Hence $f(p)\equiv\sum_{k=0}^m\binom mk(-4)^k=(-3)^{(p-1)/2}$, which by Euler's criterion and $\bigl(\frac{-3}{p}\bigr)=\bigl(\frac p3\bigr)$ is $\equiv0,1,-1$ according as $p=3$, $p\equiv1$, $p\equiv2\pmod 3$ (and $f(2)=1$).

\problem{10}
\hint{1} Get rid of the root: $(a_{n+2}-5(a_{n+1}+a_n))^2=4(6a_{n+1}^2+15a_{n+1}a_n+6a_n^2)$. Expand and look at the symmetry of the result.
\hint{2} It becomes $a_{n+2}^2+a_{n+1}^2+a_n^2=10(a_{n+2}a_{n+1}+a_{n+1}a_n+a_na_{n+2})$, symmetric in the three consecutive terms. So for $n\ge2$ both $a_{n+2}$ and $a_{n-1}$ are roots of $z^2-10(a_{n+1}+a_n)z+(a_{n+1}^2+a_n^2-10a_{n+1}a_n)=0$.
\hint{3} The sequence is increasing, so $a_{n-1}<5(a_{n+1}+a_n)\le a_{n+2}$: they are the two different roots, and Vieta gives $a_{n+2}=10(a_{n+1}+a_n)-a_{n-1}$. With $a_3=71$ computed directly, integrality follows by induction.

\problem{11}
\subpart{a}
\hint{1} Compute: $1,2,5,14,41,122,\dots$; guess $a_n=\frac{3^n+1}2$.
\hint{2} Verify by induction: $\frac{3^{n+1}+1}2\cdot\frac{3^{n-1}+1}2=3^{n-1}+\bigl(\frac{3^n+1}2\bigr)^2$ (both sides minus $\frac{3^{2n}+1}{4}$ equal $\frac{10\cdot3^{n-1}}4$). Equivalently $a_{n+1}=4a_n-3a_{n-1}$, which gives integrality directly.
\subpart{b}
\hint{1} $\frac1{a_n}=\frac2{3^n+1}<2\cdot3^{-n}$: compare with a geometric series. It converges.
\subpart{c}
\hint{1} By (a), $\frac{2^n}{a_{2^n}}=\frac{2^{n+1}}{y_n+1}$ with $y_n=3^{2^n}$, and $y_{n+1}=y_n^2$. Check the identity $\dfrac1{y+1}=\dfrac1{y-1}-\dfrac2{y^2-1}$.
\hint{2} Hence $\dfrac{2^{n+1}}{y_n+1}=\dfrac{2^{n+1}}{y_n-1}-\dfrac{2^{n+2}}{y_{n+1}-1}$, the series telescopes and its sum is $\dfrac{2}{y_0-1}=1$ (the remainder $\frac{2^{N+1}}{3^{2^N}-1}\to0$). Rational.
\end{document}
